\documentclass{article}
\usepackage[margin=1in]{geometry}
\usepackage{parskip}
\usepackage{amsmath}

\title{ISE 555 HW 2}
\author{Luke Pepin}
\date{September 30 2026}

\begin{document}
\maketitle

\section*{Question 2}
Write $(2,2) = \alpha_1(0,0) + \alpha_2(1,4) + \alpha_3(3,1)$ with $\alpha_1 + \alpha_2 + \alpha_3 = 1$ and all $\alpha_i \ge 0$. One equation for each coordinate:

$\alpha_2 + 3\alpha_3 = 2$

$4\alpha_2 + \alpha_3 = 2$

From the first equation $\alpha_3 = 2/3 - (1/3)\alpha_2$

$4\alpha_2 + 2/3 - (1/3)\alpha_2 = 2$ so $(11/3)\alpha_2 = 4/3$ and $\alpha_2 = 4/11$

$3\alpha_3 = 2 - 4/11 = 18/11$ so $\alpha_3 = 6/11$

$\alpha_1 = 1 - 4/11 - 6/11 = 1/11$

$(2,2) = (1/11)(0,0) + (4/11)(1,4) + (6/11)(3,1)$, with all weights $\ge 0$

\section*{Question 3}
\textbf{Part (A)} $f(x) = 3x^2 + 4x - 5$

$f'(x) = 6x + 4$, $f''(x) = 6 > 0$ for all $x$, so convex

\textbf{Part (B)} $f(x) = 4 - 5x + 3x^2$

$f'(x) = -5 + 6x$, $f''(x) = 6 > 0$ for all $x$, so convex

\textbf{Part (C)} $f(x) = \exp(x^2) = e^{x^2}$

Chain rule: $f'(x) = 2xe^{x^2}$

Product rule: $f''(x) = 2e^{x^2} + 4x^2e^{x^2} = (2 + 4x^2)e^{x^2}$

$e^{x^2} > 0$ and $2 + 4x^2 \ge 2$, so $f''(x) > 0$ for all $x$, so convex

\textbf{Part (D)} $f(x) = 7x - 15$

$f'(x) = 7$, $f''(x) = 0$ for all $x$. Since $0 \ge 0$ and $0 \le 0$ it is both convex and concave (a linear function is both at every $x$)

\section*{Question 4}
\textbf{Part (A)} At each point, check which constraints are active, then write the conditions on $p$.
\begin{itemize}
    \item $(0,0)$: $0 + 0 < 2$ inactive, $x_1 = 0$ active, $x_2 = 0$ active, so $p_1 \ge 0$ and $p_2 \ge 0$
    \item $(0,1)$: $0 + 1 < 2$ inactive, $x_1 = 0$ active, $x_2 = 1$ inactive, so $p_1 \ge 0$ and $p_2$ is free (any real number)
    \item $(1,1)$: $1 + 1 = 2$ active, $x_1 = 1$ inactive, $x_2 = 1$ inactive, so $p_1 + p_2 \le 0$
    \item $(0,2)$: $0 + 2 = 2$ active, $x_1 = 0$ active, $x_2 = 2$ inactive, so $p_1 + p_2 \le 0$ and $p_1 \ge 0$
\end{itemize}

\textbf{Part (B)} $f(x) = -x_1 - x_2$, so $c = (-1,-1)$ and a descent direction needs $-p_1 - p_2 < 0$, which is $p_1 + p_2 > 0$. A point can be a local minimizer if no feasible $p$ is also a descent direction.
\begin{itemize}
    \item $(0,0)$: $p = (1,1)$ is feasible and $1 + 1 > 0$, so not a local minimizer
    \item $(0,1)$: $p = (1,1)$ is feasible and $1 + 1 > 0$, so not a local minimizer
    \item $(1,1)$: needs $p_1 + p_2 \le 0$ and $p_1 + p_2 > 0$, not possible, so it can be a local minimizer
    \item $(0,2)$: needs $p_1 + p_2 \le 0$ and $p_1 + p_2 > 0$, not possible, so it can be a local minimizer
\end{itemize}

$(1,1)$ and $(0,2)$ can be local minimizers.

\section*{Question 6}
\textbf{Part (A)} At $(0,1)$ the equality and $x_1 \ge 0$ are active, $x_2 \ge 0$ is inactive.

$p_1 + p_2 = 0$ and $p_1 \ge 0$, let $p_1 = t$ so $p_2 = -t$

$p = t(1,-1)$ with $t \ge 0$

\textbf{Part (B)} At $(1,0)$ the equality and $x_2 \ge 0$ are active, $x_1 \ge 0$ is inactive.

$p_1 + p_2 = 0$ and $p_2 \ge 0$, let $p_2 = t$ so $p_1 = -t$

$p = t(-1,1)$ with $t \ge 0$

\textbf{Part (C)} At $(0.5,0.5)$ only the equality is active, $x_1 \ge 0$ and $x_2 \ge 0$ are both inactive.

$p_1 + p_2 = 0$ with no inequality restricting it, let $p_2 = t$ so $p_1 = -t$

$p = t(-1,1)$ with $t$ any real number

\section*{Question 7}
The global optimum is the single best answer out of every possible option, the local optimum an answer that no small change can improve, it's the best among its nearby options. In terms for the now 85 machine size factory, doing a global optimum search is time consuming, expensive and potentially disruptive to the whole business operation, a local optimum is highly efficient here since we already have a highly-optimized schedule from the 84 machines and we are now just seeking to improve it slightly not upend the whole thing.

\end{document}
